\begin{figure}[H]
    \centering

    \begin{tikzpicture}[
        font=\small,
        box/.style={
            draw,
            rounded corners,
            minimum width=3.4cm,
            minimum height=1cm,
            align=center,
            fill=gray!10
        },
        store/.style={
            draw,
            cylinder,
            shape border rotate=90,
            aspect=0.25,
            minimum width=3.2cm,
            minimum height=1.2cm,
            align=center,
            fill=gray!20
        },
        arrow/.style={
            -Latex,
            thick
        },
        dashedarrow/.style={
            -Latex,
            thick,
            dashed
        }
    ]

\node[box] (broker) at (0,3) {Broker de Eventos\\Kafka / RabbitMQ};

\node[box] (consumer) at (0,1) {Consumidor\\Idempotente};

\node[store] (processed) at (-4,-1) {Eventos\\Procesados};

\node[box] (business) at (4,-1) {Lógica de\\Negocio};

\draw[arrow] (broker) -- node[right]{Evento} (consumer);

\draw[arrow] (consumer) -- node[above]{1. Verificar ID} (processed);

\draw[arrow] (consumer) -- node[above]{2. Procesar si no existe} (business);

\draw[dashedarrow] (processed) -- node[below]{Evento duplicado} (consumer);

    \end{tikzpicture}

    \caption{Control de eventos duplicados mediante consumidor idempotente.}
    \label{fig:idempotencia_eventos}

\end{figure}